\chaptermark{Bidirectional cardinal closure of the analytic continuum HMT}
\label{chap:83-19}

This chapter does not propose an alternative generator of the continuum. Its
causal premise is the unique block already constructed
\[
 \begin{aligned}
 \mathrm{APP}&\longrightarrow\mathrm{TRIT}\longrightarrow\mathrm{TPK}
 \longrightarrow\text{unique enriched state}\\
 &\longrightarrow\text{common state, tree and genealogical measure}\\
 &\longrightarrow\text{correlative intrinsic operations}\\
 &\longrightarrow\text{natural assembly, terminality and limit}
 \longrightarrow\mathcal C_{\rm cont}^{\rm disc}.
 \end{aligned}
\]
The five consubstantial constructions developed in the preceding chapter emerged correlatively over that same state before the real publication, without purporting to exhaust all its properties. What follows is a subsequent presentation of the quotient, its temporal memory, and its power; it does not retrospectively select any of those constructions.

\section{The complete TRIT cell and the bidirectional temporal extension}

The temporal sign does not separate from the state it orients. The local cell of
departure is
\begin{equation}
 c=(i,j;R^+,Q^+,R^\times,Q^\times;\tau),
 \qquad \tau\in\{-1,0,+1\},
 \label{p12-83-c19-s1:hmtch:eq-celula-trit-completa}
\end{equation}
where
\[
 R^+=\operatorname{dr}_9(i+j),\qquad
 Q^+=\frac{i+j-R^+}{9},
\]
\[
 R^\times=\operatorname{dr}_9(ij),\qquad
 Q^\times=\frac{ij-R^\times}{9}.
\]
The residues constitute the visible reading and the quotients preserve the
memory necessary for inverting the transport. The temporal trit acts on
this integer cell:
\begin{equation}
\begin{aligned}
 +1&\longmapsto \text{return, retained memory and past},\\
 0&\longmapsto \text{evaluation threshold and present},\\
 -1&\longmapsto \text{bidirectional opening and future}.
\end{aligned}
\label{p12-83-c19-s1:hmtch:eq-lector-temporal-trit}
\end{equation}
Thus, two appearances of the same sign are not the same state if their
quotients, sheet, carry, or genealogical record differ.

Let \(X\) be a compact chart of the surviving APP--TRIT--TPK limit and let
\(F:X\to X\) be its complete, continuous and surjective update. The surjectivity of the truncations and the surviving prolongation allow us to form the natural extension
\begin{equation}
 X_F^{\leftrightarrow}
 :=\left\{(x_k)_{k\in\mathbb Z}\in X^{\mathbb Z}:
 F(x_k)=x_{k+1}\ \text{for every }k\in\mathbb Z\right\}.
 \label{p12-83-c19-s1:hmtch:eq-extension-natural-two-way}
\end{equation}
The displacement
\[
 \sigma((x_k)_{k\in\mathbb Z})=(x_{k+1})_{k\in\mathbb Z}
\]
it no longer loses the inverse direction.

\begin{theorem}[Natural extension two-way of the enriched state]
\label{p12-83-c19-s1:hmtch:thm-extension-natural-two-way}
The space \(X_F^{\leftrightarrow}\) is compact and

\[
 \sigma:X_F^{\leftrightarrow}\longrightarrow X_F^{\leftrightarrow}
\]
is a homeomorphism, with
\[
 \sigma^{-1}((x_k))=(x_{k-1}).
\]
Therefore, outward and return maps are inverse continuous morphisms on the
enriched object. The apparent discontinuity arises only if the coordinates
that make the history invertible are first forgotten.
\end{theorem}

\begin{proof}
The product \(X^{\mathbb Z}\) is compact. The condition
\(F(x_k)=x_{k+1}\) is closed becau \(F\) is continuous; by tanto
\(X_F^{\leftrightarrow}\) is compact. The displazamiento and the
displazamiento inver are continuous in the topology product and priservan
the equations of compatibility. Their compositions are the identity.
\end{proof}

The material inversion of the route resides in the groupoid completion of
the transport. Let \(C_M=-\operatorname{rev}\) be memory conjugation. Its
contravariant action is expressed by the reversal square
\[
 F(x_k)=x_{k+1}
 \quad\Longrightarrow\quad
 F(C_Mx_{k+1})=C_Mx_k.
\]
On the natural extension it is defined
\begin{equation}
 (\Theta u)_k:=C_M(x_{-k}),
 \qquad u=(x_k)_{k\in\mathbb Z}.
 \label{p12-83-c19-s1:hmtch:eq-definicion-espejo-two-way}
\end{equation}
The reversion identity guarantees that \(\Theta u\) again satisfies
the compatibility equations. Therefore,
\begin{equation}
 \Theta^2=I,
 \qquad
 \Theta\sigma\Theta=\sigma^{-1}.
 \label{p12-83-c19-s1:hmtch:eq-espejo-inversion-two-way}
\end{equation}
This identity does not identify the ternary group with the binary sign: the
TRIT orients return, threshold, and opening, whereas the involution \(C_M\)
implements the conjugation of the two temporal orientations.

For \(n\geq0\), let us denote by \(S_n^{\leftrightarrow}\) the finite set
of admissible temporal windows
\((x_{-n},\ldots,x_0,\ldots,x_n)\), with each state further truncated to
genealogical depth \(n\). The simultaneous restrictions of time and
depth form an
inverse system and
\begin{equation}
\mathscr O_{\leftrightarrow}^{\rm HMT}
 :=\varprojlim_n S_n^{\leftrightarrow}
 \cong X_F^{\leftrightarrow}.
 \label{p12-83-c19-s1:hmtch:eq-biordinal-genealogico}
\end{equation}
Each history has temporal support of type
\(\omega^*+1+\omega\), but also carries the cell
\eqref{p12-83-c19-s1:hmtch:eq-celula-trit-completa}, the route and the genealogical register. The map sending an orbit to all its centered windows is a homeomorphism and intertwines \(\sigma\) with the genealogical succession. This is the genealogical biordinal, not the von Neumann ordinal obtained after forgetting the labels.

The preceding object is a compact chart of the fractional part. The global
object publishing the entire line is the countable atlas
\begin{equation}
 \mathscr O_{\mathbb R,\leftrightarrow}^{\rm HMT}
 :=\mathbb Z\times\mathscr O_{\leftrightarrow}^{\rm HMT}.
 \label{p12-83-c19-s1:hmtch:eq-biordinal-global}
\end{equation}
The coordinate \(\mathbb Z\) merely translates the chart; it does not erase the bidirectional
genealogical ledger.

\section{Unitary calendar, modular memory and stable mode}

The temporal structure does not rest solely on a metaphor of return.
On
\(\mathcal H_{\rm clk}=\mathbb C^{108}\), let
\(\zeta=\exp(2\pi i/108)\),
\[
 U_{\rm clk}|j\rangle=|j+1\rangle,
 \qquad
 |\widetilde k\rangle
 =\frac1{\sqrt{108}}\sum_{j=0}^{107}\zeta^{jk}|j\rangle,
\]
and let us define
\begin{equation}
 H_{\rm clk}
 =\hbar_{\rm int}\omega_0
 \sum_{k=0}^{107}k\,|\widetilde k\rangle\langle\widetilde k|,
 \qquad
 \omega_0=\frac{2\pi}{108t_0}.
 \label{p12-83-c19-s1:hmtch:eq-hamiltoniano-reloj}
\end{equation}
Then
\begin{equation}
 U_{\rm step}
 =\exp\!\left(-\frac{i}{\hbar_{\rm int}}H_{\rm clk}t_0\right)
 =U_{\rm clk},
 \qquad U_{\rm step}^{108}=I,
 \label{p12-83-c19-s1:hmtch:eq-propagador-reloj}
\end{equation}
and, for \(Z_{\rm clk}|j\rangle=\zeta^j|j\rangle\),
\[
 |\psi_n\rangle:=U_{\rm step}^n|j_0\rangle,
\]
\begin{equation}
 \langle Z_{\rm clk}\rangle_n
 :=\langle\psi_n|Z_{\rm clk}|\psi_n\rangle
 =\zeta^{j_0+n}
 \label{p12-83-c19-s1:hmtch:eq-respuesta-primitiva-108}
\end{equation}
has primitive period \(108\).

In the modular boundary sector
\(\mathcal H_A=(\mathbb C^3)^{\otimes36}\), let \(\rho_A>0\) be precisely
the faithful state characterized by the following decomposition, and let
\(\Delta_4^{\rm mod}\) be the declared modular jump. The generator preserving
the provenance of the two boundary channels is
\begin{equation}
 K_{\rm HHM}^{(A)}=-\log\rho_A
 =\kappa_A I+\Delta_4^{\rm mod}(90N_{90}+120N_{120}).
 \label{p12-83-c19-s1:hmtch:eq-hamiltoniano-modular-memoria}
\end{equation}
If
\[
 N=N_{90}+N_{120},
 \qquad
 D=90N_{90}+120N_{120},
\]
then
\begin{equation}
 N_{90}=4N-\frac{D}{30},
 \qquad
 N_{120}=\frac{D}{30}-3N.
 \label{p12-83-c19-s1:hmtch:eq-reconstruccion-canales-modulares}
\end{equation}
Thus, the total observable and the modular Hamiltonian jointly reconstruct
the two provenance memories. On the other hand, the cylindrical Perron
dynamics preserves the mode
\begin{equation}
 V_n=\varphi^n,
 \qquad
 M_n=\varphi^{-2n}=V_n^{-2},
 \qquad
 M_n=\left(\frac{a_n}{a_0}\right)^{-6}\qquad(D=3).
 \label{p12-83-c19-s1:hmtch:eq-modo-estable-memoria}
\end{equation}

\begin{theorem}[Exact finite composition of the CT108 calendar and modular memory \(90/120\)]
\label{p12-83-c19-s1:hmtch:thm-composicion-calendario-memoria}
The operators
\eqref{p12-83-c19-s1:hmtch:eq-hamiltoniano-reloj} and
\eqref{p12-83-c19-s1:hmtch:eq-hamiltoniano-modular-memoria} are self-adjoint on their
respective spaces. After an energy scale \(E_\partial>0\) is fixed, the
product operator
\begin{equation}
 H_{\rm tw}
 :=H_{\rm clk}\otimes I
 +I\otimes E_\partial(K_{\rm HHM}^{(A)}-\kappa_A I)
 \label{p12-83-c19-s1:hmtch:eq-hamiltoniano-two-way-total}
\end{equation}
is self-adjoint, and its two summands strongly commute. For states whose
clock marginal is \(|j_0\rangle\langle j_0|\), the observable
\(Z_{\rm clk}\otimes I\) preserves its primitive response of period \(108\).
Simultaneously, \(K_{\rm HHM}^{(A)}\), \(N_{90}\), and \(N_{120}\) are
conserved observables of the modular factor. The pair \((N,D)\), with
\(N=N_{90}+N_{120}\), reconstructs \((N_{90},N_{120})\) exactly by means of
\eqref{p12-83-c19-s1:hmtch:eq-reconstruccion-canales-modulares}.
\end{theorem}

\begin{proof}
The first operator admits a finite real spectral decomposition, and the
second is the logarithm of a faithful state; both are self-adjoint. They act on
distinct tensor factors and, therefore, their spectral resolutions
commute. The exponential of their sum factorizes. The primitivity of the response
follows from the primitivity of \(\zeta\); channel reconstruction is the
inverse of the matrix
\(\left(\begin{smallmatrix}1&1\\90&120\end{smallmatrix}\right)\), de
determinante \(30\).
\end{proof}

\begin{corollary}[Stable mode of the cylindrical memory]
\label{p12-83-c19-s1:hmtch:cor-modo-estable-memoria}
In the Perron sector of closed paths,
\(M_n=\varphi^{-2n}=V_n^{-2}\) is the exact and stable spectral mode given
by \eqref{p12-83-c19-s1:hmtch:eq-modo-estable-memoria}.
\end{corollary}

The CT108 calendar, the modular memory and the Perron mode are three readeres
tipados of the same genealogical antecedent; are not the same operator. together with
the natural extension and the espejo, form the internal temporal record
\[
 \mathcal T_{\rm HMT}
 :=\bigl(X_F^{\leftrightarrow},\sigma,\Theta;
 H_{\rm clk},K_{\rm HHM}^{(A)};M_\bullet\bigr).
\]
This is the exact content of the HMT temporal closure: two inverse
continuous directions, a primitive clock, recoverable memory, and a stable mode. The
complete TPK state is not replaced by the cursor \(108\). An additional
identification of the three readers with a single perturbatively stable physical
dynamics would require a common intertwiner; that identification is not used in
the genealogical closure or in the reconstruction of the continuum.

\section{Analytic publications of HMT histories}

Let \(H\) be a Polish space of histories: a closed inverse limit of
finite levels or a countable atlas of such limits. The topology is
generated by prefix cylinders. This choice preserves the genealogy and
permits the descriptive operations that are admitted to be typed exactly.

\begin{definition}[Analytic HMT Publication]
\label{p12-83-c19-s1:hmtch:def-publicacion-analitica}
An analytic publication is a set \(A\subseteq\mathbb R\) for which
there exists a Borel set
\[
 C\subseteq H\times\mathbb N^{\mathbb N}
\]
and to Borel reader \(F:H\to\mathbb R\) such that
\[
 A=\{F(x):\exists z\in\mathbb N^{\mathbb N},\ (x,z)\in C\}.
\]
The witness \(z\) may encode to route, to prolongation, to memory or
to certification countable. The readers continuous are cases particular.
\end{definition}

This class contains cylinders, Borel predicates on histories, countable
operations, images under continuous readers, and publications asserting the
existence of a real witness. By definition, it does not contain the full power set
\(\mathcal P(\mathbb R)\): uncountable unions, arbitrary selectors, nonsmooth quotients, and
recursion of length \(\omega_1\) require additional data.

\begin{theorem}[Perfect Duality of HMT Publications]
\label{p12-83-c19-s1:hmtch:thm-dicotomia-perfecta}
Every analytical HMT publication \(A\subseteq\mathbb R\) satisfies exactly
one of the two alternatives:
\[
 |A|\leq\aleph_0
 \qquad\text{or}\qquad
 |A|=2^{\aleph_0}.
\]
In the second alternative there exists a continuous injection.
\[
 j:2^{\mathbb N}\hookrightarrow A.
\]
therefore, every uncountable analytic publication contains one copia of the
Cantor space and, in particular, a perfect subset.
\end{theorem}

\begin{proof}
The existential projection of a Borel subset of a product of Polish spaces is analytic, and a Borel image of a standard Borel space is again analytic. Therefore, \(A\) is an analytic subset of
\(\mathbb R\).

To make the uncountable alternative explicit, represent \(A\) as
the projection of a closed set
\[
 D\subseteq\mathbb R\times\mathbb N^{\mathbb N}.
\]
For each finite word \(s\in\mathbb N^{<\mathbb N}\), let
\[
 A_s:=\{x:\exists z\supseteq s,\ (x,z)\in D\}.
\]
If \(A_s\cap K\) is uncountable for a certain closed interval \(K\), the set of its points that are not condensation points is countable, because \(\mathbb R\) has a countable basis. Two distinct condensation points and two disjoint closed intervals within \(K\) are chosen, and for each interval a strict prolongation of \(s\) whose projection remains uncountable is chosen. The two prolongations need not be incompatible: separation of the branches is guaranteed by the intervals. The identity
\[
 A_s=\bigcup_{n\in\mathbb N}A_{s^\frown n}
\]
guarantees that, in each chosen neighborhood, some extension preserves the non
countability.

Repitiendo the procedure are obtained, for each
\(u\in2^{<\mathbb N}\), a interval \(K_u\) and a prefix \(s_u\) such that
\[
 K_{u0}\cap K_{u1}=\varnothing,
 \qquad
 K_{ui}\subseteq K_u,
 \qquad
 \operatorname{diam}K_u\leq2^{-|u|},
\]
and \(s_{ui}\) strictly prolongs to \(s_u\). For
\(\beta\in2^{\mathbb N}\), the intervals
\(K_{\beta|n}\) determine a unique point \(j(\beta)\), and the prefixes
\(s_{\beta|n}\) determine a witness \(z_\beta\). The closure of \(D\)
implies
\[
 (j(\beta),z_\beta)\in D,
\]
so that \(j(\beta)\in A\). The disjoint sibling intervals prove
injectivity, and diameter control proves the continuity of \(j\).

Thus,
\[
 2^{\aleph_0}\leq |A|\leq|\mathbb R|=2^{\aleph_0},
\]
which closes the second alternative. If the process does not start from an
uncountable set, the first is obtained.
\end{proof}

\begin{corollary}[Analytic HMT Continuum Hypothesis]
\label{p12-83-c19-s1:hmtch:cor-ch-analitica-hmt}
Let
\begin{equation}
 \aleph_{1,\mathrm{HMT}}^{\rm an}
 :=\min\left\{|A|:
 A\text{ is an uncountable analytic HMT publication}\right\}.
 \label{p12-83-c19-s1:hmtch:eq-aleph1-analitico-hmt}
\end{equation}
This symbol denotes the initial cardinal relative to the class of HMT analytic
publications; it does not denote the extensional ordinal \(\omega_1\).
Then
\begin{equation}
 \boxed{
 \aleph_{1,\mathrm{HMT}}^{\rm an}
 =2^{\aleph_0}
 =|\mathbb R|
 =\left|\mathscr O_{\mathbb R,\leftrightarrow}^{\rm HMT}/
 \!\sim_{\mathfrak R}\right|.}
 \label{p12-83-c19-s1:hmtch:eq-ch-analitica-hmt}
\end{equation}
That is, there is no intermediate cardinal between the countable and the continuum
within the analytic publications produced by HMT readers.
\end{corollary}

\begin{proof}
The set of finite depth windows
\[
 \mathscr O_{<\omega}^{\rm HMT}
 :=\bigcup_{n<\omega}S_n^{\leftrightarrow}
\]
is countable, since it is a countable union of finite sets. The complete layer
surjects onto \(\mathbb R\) by
\eqref{p12-83-c18-s2:hmtch:eq-homeomorfismo-two-way-real} and is contained in a countable product
of finite alphabets; it therefore has cardinality
\(2^{\aleph_0}\). The
\cref{p12-83-c19-s1:hmtch:thm-dicotomia-perfecta} excludes any strictly intermediate analytic
size.
\end{proof}

\begin{corollary}[Dichotomy of the fibers]
\label{p12-83-c19-s1:hmtch:cor-dicotomia-fibras}
Let \(F:H\to Y\) to reader Borel between spaces Polish and let
\(y\in Y\). If the domain of histories admitted is analytic, the fiber
\[
 \{x:F(x)=y\}
\]
is countable or has cardinality \(2^{\aleph_0}\).
\end{corollary}

\begin{proof}
The fiber is the intersection of the analytic domain with the Borel preimage of the singleton \(\{y\}\); it is therefore analytic. We apply the
\cref{p12-83-c19-s1:hmtch:thm-dicotomia-perfecta}.
\end{proof}

\section{Comparison between the genealogical cardinal closure and its extensional ordinal reduction}

Let \(\operatorname{WO}\subseteq2^{\mathbb N}\) be the set of real
codes of countable well-orders, allowing the field to be a
subset of \(\mathbb N\). The map
\[
 \operatorname{ot}:\operatorname{WO}\longrightarrow\omega_1
\]
assigns to each code its order type and is surjective.

\begin{proposition}[Ordinal injection into the line]
\label{p12-83-c19-s1:hmtch:prop-omega1-en-reales}
In ZFC there exists an injection
\[
 \iota:\omega_1\hookrightarrow\mathbb R.
\]
A binary subsystem of the history space HMT can be raised.
\end{proposition}

\begin{proof}
By choice, the surjection \(\operatorname{ot}\) admits a section
\(s:\omega_1\to\operatorname{WO}\). Equivalently, a well-order
of \(\operatorname{WO}\) is fixed and the first code in each fiber is taken.
The Cantor map
\[
 c(a):=\sum_{n\geq0}\frac{2a(n)}{3^{n+1}}
\]
is injective from \(2^{\mathbb N}\) into \([0,1]\). Therefore
\[
 \iota:=c\circ s
\]
is the injection sought. The embedding of the binary tape in the cinta
ternary and the Hensel inverse \(\Psi^{-1}\) proporcionan the lift to the
genealogical space.
\end{proof}

The choice of \(s\) is not natural with respect to refinements and is not
produced by a finite prefix. This proposition shows that the arrow
\(\omega_1\hookrightarrow\mathbb R\) is not the content of the classical
hypothesis of the continuum. The inverse direction belongs to the bare extensional
ordinal. In HMT, by contrast, the natural arrow already constructed is the
homeomorphism \eqref{p12-83-c18-s2:hmtch:eq-homeomorfismo-two-way-real}, whose domain
preserves the two temporal directions.

\section{Cardinal identity of the HMT analytic continuum: \(\aleph_{1,\mathrm{HMT}}^{\mathrm{an}}=2^{\aleph_0}\) and comparison with the classical continuum hypothesis}

Let
\[
 Q_{\mathrm{HMT}}:=X_{\mathrm{raw}}/\!\sim_P.
\]
By the \cref{p12-83-c18-s2:hmtch:thm-cociente-real},
\(Q_{\mathrm{HMT}}\cong\mathbb R\).

\begin{theorem}[Ordinal equivalences of the HMT quotient]
\label{p12-83-c19-s1:hmtch:thm-equivalencias-ch}
In ZFC the following are equivalent:
\begin{enumerate}
 \item \(2^{\aleph_0}=\aleph_1\);
 \item there exists a injection \(j:Q_{\mathrm{HMT}}\hookrightarrow\omega_1\);
 \item there exists to map \(\rho:X_{\mathrm{raw}}\to\omega_1\) such that
 \[
  \rho(x)=\rho(y)
  \quad\Longleftrightarrow\quad
  P(x)=P(y);
 \]
 \item \(Q_{\mathrm{HMT}}\) admits an uncountable well-order whose proper
 initial segments are all countable.
\end{enumerate}
\end{theorem}

\begin{proof}
If the first assertion holds, a bijection \(Q_{\mathrm{HMT}}\cong\mathbb R\cong\omega_1\) provides the second and transports the order of \(\omega_1\), which yields the fourth. The second yields the third by taking \(\rho=j\circ[P]\), where \([P]\) sends each history to its evaluation class. Conversely, the biconditional condition in the third causes \(\rho\) to factor through an injection \(Q_{\mathrm{HMT}}\hookrightarrow\omega_1\).

The second assertion implies
\(2^{\aleph_0}\leq\aleph_1\). The
\cref{p12-83-c19-s1:hmtch:prop-omega1-en-reales} gives
\(\aleph_1\leq2^{\aleph_0}\); Cantor--Bernstein produces the equality. Finally,
a well-order such as that of the fourth assertion cannot contain an
element in position \(\omega_1\), because its initial segment would be uncountable.
Its order type is therefore at most \(\omega_1\), and, since
the domain is uncountable, it is exactly \(\omega_1\). This recovers the
second assertion.
\end{proof}

The theorem localizes the target without introducing it as a premise. A function
\(\rho\) with that property would not be a small supplement to the analytic
dichotomy: it would be a complete formulation of the continuum
hypothesis itself.

\section{Loss of ancestral memory under reduction to the ordinal codomain}

The forgetful functor that removes trit, mirror, quotients, route, and memory sends
each biordinal to a bare ordered support. The following results
control that forgetful functor; they do not deny the continuous HMT morphism already
constructed between the two-way quotient and \(\mathbb R\).

\begin{theorem}[Limit of the readers toward \(\omega_1\) with the order topology]
\label{p12-83-c19-s1:hmtch:thm-no-go-continuo}
There does not exist a continuous application
\[
 \rho:X_{\mathrm{raw}}\longrightarrow\omega_1
\]
that exactly separates the fibers of \(P\). In particular, after
applying the forgetful functor, no cylindrical reader toward the bare
topological ordinal can demonstrate the second condition of
\cref{p12-83-c19-s1:hmtch:thm-equivalencias-ch}.
\end{theorem}

\begin{proof}
Every chart \(\{m\}\times\mathbb Z_3^6\) is compact. The continuous image of
a compact space in \(\omega_1\), endowed with the order topology, is compact and
therefore bounded by some countable ordinal. Its image is countable.
The union over \(m\in\mathbb Z\) of those images is a countable union of
countable sets and remains countable. A function separating all
fibers of \(P\) would instead have an image of cardinality
\(|\mathbb R|\), a contradiction.
\end{proof}

\begin{theorem}[Absence of good analytic extensional order]
\label{p12-83-c19-s1:hmtch:thm-no-go-analitico}
There is no to good order analytic, and by tanto nor Borel, of
\(\mathbb R\).
\end{theorem}

\begin{proof}
Suppose that \(\prec\subseteq\mathbb R^2\) is analytic and defines a strict
total linear order. Then
\[
 (x,y)\notin\prec
 \quad\Longleftrightarrow\quad
 x=y\ \text{or}\ y\prec x.
\]
The complement of \(\prec\) is also analytic; by Suslin separation, \(\prec\) is Borel. Every well-founded Borel relation on a Polish space has rank bounded by a countable ordinal. In a linear well-order, distinct elements have distinct ranks, so the domain would be countable. This contradicts the uncountability of \(\mathbb R\).
\end{proof}

The two controls are independent of whether the classical continuum hypothesis is
true or false. If a set-theoretic bijection
\(\omega_1\leftrightarrow\mathbb R\) exists, it cannot be continuous for the
order topology of \(\omega_1\) or analytic as a well-order. This does not make
the HMT morphism discontinuous: it shows that bare \(\omega_1\) is another
codomain and that forgetting the temporal direction destroys precisely the
structure that ensured continuity.

\section{Ordinal closure of the HMT genealogical codomain}

The preceding checks describe the functor that forgets genealogy. HMT
positively constructs its ordinal closure while preserving it. Let
\(h\in2^{\mathbb N}\) be the generated code of the language, finite tables,
transitions and APP--TRIT--TPK readers. The canonical genealogical
codomain is
\[
 M_h:=L[h].
\]
This envelope preserves the generator as a parameter, admits the canonical
well-order of the constructible hierarchy and is the codomain produced by the
construction of the continuum in \eqref{eq:frontal-doble-proyeccion}.

\begin{theorem}[Classical continuum hypothesis in the HMT genealogical codomain]
\label{p12-83-c19-s1:hmtch:thm-ch-relativa}
The model \(M_h=L[h]\) satisfies
\[
 M_h\models 2^{\aleph_0}=\aleph_1.
\]
There exists in \(M_h\) a definable bijection
\[
 b_h:\omega_1^{M_h}\xrightarrow{\ \sim\ }\mathbb R^{M_h}.
\]
Since \(M_h\) is the generated codomain, \(\mathbb R^{M_h}\) is the entirety of
the line published by the HMT construction. Let \(P_h\) be the HMT evaluation of
the histories belonging to \(M_h\). Then
\[
 \rho_h:=b_h^{-1}\circ P_h
\]
satisfies
\[
 \rho_h(x)=\rho_h(y)
 \quad\Longleftrightarrow\quad
 P_h(x)=P_h(y).
\]
\end{theorem}

\begin{proof}
Every real number in \(L[h]\) appears at some level
\(L_\alpha[h]\) with \(\alpha<\omega_1^{L[h]}\). For each ordinal
\(\alpha<\omega_1^{L[h]}\), the level \(L_\alpha[h]\) is countable within
the model. Therefore,
\[
 |\mathbb R^{L[h]}|\leq\aleph_1^{L[h]}.
\]
The model line is uncountable, so minimality of \(\aleph_1\)
gives the reverse inequality and hence equality. The canonical well-ordering
of \(L[h]\) selects, among all the bijections whose existence has just
been proved, a first definable bijection \(b_h\).

The evaluation \(P_h\) is definable from \(h\), so that its value
belongs to \(\mathbb R^{M_h}\). The formula of \(\rho_h\) is well typed.
Since \(b_h\) is bijective, two histories receive the same ordinal if and only
if they publish the same real.
\end{proof}

\begin{remark}[Identification of the generated codomain]
\label{p12-83-c19-s1:hmtch:rem-alcance-relativo}
\cref{p12-83-c19-s1:hmtch:thm-ch-relativa} effectively constructs the
ordinal arrow in the genealogically generated universe. The equality
\(\mathbb R=\mathbb R^{L[h]}\) is not postulated: within the theorem's domain
it expresses that \(L[h]\) is precisely the codomain produced by the HMT
continuum, whose Archimedean coverage is proved in
\cref{p12-83-c18-s1:prop:agotamiento-recta-real-u024,p12-83-c18-s2:hmtch:thm-cociente-real}.
Consequently, \(\rho_h\) reaches all the reals of the generated line and
the bijection \(b_h\) proves the classical continuum hypothesis with the
ordinary cardinals of \(M_h\). A forcing extension belongs to another
domain after forgetting or replacing the genealogical selector \(h\); it describes
the independence of the bare language of ZFC and does not alter the conclusion of the
HMT system.
\end{remark}

\section{Compatibility among HMT cardinal closure, the set-theoretic formulation of the continuum, and bidirectional genealogical memory}

The preceding results articulate five strata of the same closure:
\begin{enumerate}
 \item the genealogical quotient reconstructs \(\mathbb R\) and has cardinality
       \(2^{\aleph_0}\);
 \item the natural extension two-way and the genealogical biordinal are
       homeomorphic, with outgoing and return continuous;
 \item the clock Hamiltonian, the modular Hamiltonian, and the stable mode
       preserve calendar, provenance, and memory without reducing the TPK to the
       visible cursor;
 \item every analytic HMT publication is countable or has the cardinal of the
       continuum, so that
       \(\aleph_{1,\mathrm{HMT}}^{\rm an}=2^{\aleph_0}\);
 \item the constructible codomain \(L[h]\) has a complete ordinal enumeration
       of all generated reals and satisfies the classical
       continuum hypothesis with its ordinary cardinals.
\end{enumerate}
The first three construct the HMT object and its dynamics; the fourth closes the
cardinality in the domain of publications; the fifth supplies its ordinal
coordinate without changing the object. The identity
\(\aleph_{1,\mathrm{HMT}}^{\rm an}=2^{\aleph_0}\), together with the bijection
\(b_h\), constitutes the HMT proof of the extensional statement about the
generated line. The selector \(h\) preserves the genealogy that the bare
language of ZFC omits.

\begin{proposition}[Falsifiers]
\label{p12-83-c19-s1:hmtch:prop-falsadores}
Each assertion has an independent negative control:
\begin{enumerate}
 \item a real without preimage under \(P\) would falsify the reconstruction of the
       quotient;
 \item an intermediate state not recoverable from the complete returns
       would falsify the factorization
       \eqref{p12-83-c17-s7:hmtch:eq-factorizacion-excepcional-retorno};
 \item the loss of the primitive period, of the reconstruction \(90/120\) or
       of the mode \(\varphi^{-2n}\) would falsify one of the three readers of the
       typed HMT temporal record;
 \item to publication analytic not countable without copia of Cantor would falsify
       the dichotomy perfect;
 \item a Borel well-order of an uncountable Polish space would refute the
       analytic no-go theorem;
 \item one unbounded continuous image of one compact chart in \(\omega_1\)
       falsaria the topological not-go;
 \item the existence, in the generated codomain \(M_h\), of a real without
       an ordinal preimage under \(b_h\), or of a branch published outside
       \(L[h]\), would falsify the genealogical identification and the classical closure
       of the continuum. An extension obtained after changing the selector
       \(h\) does not belong to the fixed domain and does not constitute that falsifier.
\end{enumerate}
\end{proposition}

The structural conclusion is precise. HMT constructs a bidirectional temporal
object, publishes it Archimedeanly as \(\mathbb R\), preserves exceptional
incidence in parallel, and proves that its first uncountable analytic layer
has cardinality exactly \(2^{\aleph_0}\). The bijection
\[
 \mathscr O_{\mathbb R,\leftrightarrow}^{\rm HMT}/\!
 \sim_{\mathfrak R}
 \xrightarrow{\ \overline{\mathfrak R}\ }\mathbb R
\]
is natural and continuous. The classical ordinal \(\omega_1^{M_h}\) is the well-order
coordinate of all reals in the generated codomain. The functor that
forgets \(h\), time, mirror and memory explains the subsequent independence in
bare ZFC; it does not modify the constructed bijection or the classical closure
proved in \(M_h\).
